\section{Introduction}
\label{sec:intro}

A personalised physiological model, a \emph{digital twin}, is a simulator whose parameters are fitted to one
person's sensor data. The fitted values are then read as physiology: an insulin sensitivity, a gastric
emptying rate, a carbohydrate absorption rate. They are shown to the person, compared across people, and
used to explain or to advise. The evidence usually offered that such a twin is trustworthy is that it
predicts held-out data. This paper is about why that evidence is not enough.

A parameter can be fitted without being identified. If two parameters enter the observation only through
a combination, the data fix the combination and leave the split arbitrary; an optimizer then returns a
split anyway, it predicts held-out data perfectly well, and the number reported as physiology is an
accident of the starting point and the parameter range. Identifiability analysis exists to separate what
the data determine from what the fitting procedure supplies
\citep{bellman1970structural,cobelli1980parameter,walter1997identification}. For glucose--insulin models
it has been known to be delicate since the minimal model \citep{bergman1979,pillonetto2002minimal}. What is
new is the setting: free-living postprandial continuous glucose monitoring (CGM) with logged meals, and
simulators that are differentiable end to end, so that exact gradients, Fisher information and profile
likelihoods are available for every person rather than for a population \citep{kidger2021diffrax,
bradbury2018jax}.

We study one concrete and widely used observation, the baseline-subtracted incremental area under the
glucose curve (iAUC) of a meal, and a Bergman-type simulator with a two-compartment gut. The questions are
specific. Which of the three personalised parameters, insulin sensitivity $\SI$, gastric emptying rate
$\ke$ and carbohydrate absorption rate $\ka$, do the data determine? Is the answer a property of the
observable, of the optimizer, of the noise, or of the model being wrong? And would held-out prediction have
told us?

\paragraph{What we do.}
We combine the full (not diagonal) Fisher information matrix with its eigenvectors and the Schur
complement of the timing block, profile likelihood intervals with explicit definitions of bounded,
one-sided and flat, a generic-rank test of the stacked sensitivity matrix, and symbolic structural
identifiability. We add two instruments that a real dataset cannot provide: a \emph{replica}, in which the
model is true and the noise and carbohydrate-logging error are the subject's own, and a synthetic study
with randomly drawn true parameters, so that recovery can be read against where in the parameter range the
truth lies. We replicate on three cohorts, run the same questions on a second model class, repeat the
prediction comparison on a second cohort, and check the optimizer that produced every estimate. The
analysis plan was fixed in version control before the analyses were run; it was amended
\resAmendmentCount\ times, each amendment is listed with its reason (Appendix~\ref{app:register}), and
hypotheses declared after results were seen are marked as such. Failures are reported as failures.

\paragraph{What we find.}
\begin{enumerate}
\item \emph{$\ke$ and $\ka$ are not separable from the area, and the reason is exact.} Their exchange is
  an exact symmetry of the empty-gut model (Section~\ref{sec:theory}); the weakest Fisher direction lies
  along the exchange direction in every subject of three cohorts; and no profile interval for either
  rate is bounded under iAUC, at any of three parameter ranges and out to four times the upper bound
  (Section~\ref{sec:timing}).
\item \emph{$\SI$ is only conditionally recoverable.} Its interval is bounded for a minority of subjects at
  the primary range, the fraction rises with the range, and in a replica where the model is exactly true
  it remains low (Section~\ref{sec:si}). The weakness is largely a property of the observable and the
  design, not only of model error.
\item \emph{What is recoverable is a combination.} The sum $\tau_1=1/\ke+1/\ka$ and the product
  $p=1/(\ke\ka)$ of reciprocal rates are bounded for most subjects under the full glucose trace, and
  mostly not under the area, with or without its timing centroid (Section~\ref{sec:more}).
\item \emph{Prediction does not see any of this.} Gradient fitting and exhaustive grid search, and
  three parameters and one, give the same held-out iAUC error in two cohorts; a full-trace objective does
  not lower it (Section~\ref{sec:prediction}).
\item \emph{The same questions on another model class, and the optimizer, do not change the picture}
  (Sections~\ref{sec:dm} and~\ref{sec:more}).
\end{enumerate}

\begin{table}[t]
\centering\small
\begin{tabularx}{\linewidth}{Xll}
\toprule
claim & evidence & scope \\
\midrule
$\ke$ and $\ka$ are not separable from the area & Prop.~\ref{prop:exchange}; Sec.~\ref{sec:timing}; H4, H5, H15 & three cohorts; iAUC \\
$\SI$ is only conditionally recoverable & Sec.~\ref{sec:si}; H5, H13, H17, H19 & CGMacros; replica and synthetic \\
The combination $(\tau_1,p)$ is recoverable from the trace & Sec.~\ref{sec:more}; H11, H18, H21 & CGMacros; one model family \\
Prediction does not see any of this & Sec.~\ref{sec:prediction}; H7--H9, H20 & CGMacros and Shanghai \\
The picture holds for a second gut model & Sec.~\ref{sec:dm}; H22 & CGMacros; iAUC only \\
Clinical correlations of a fitted parameter need controls & Sec.~\ref{sec:leakage} & Shanghai \\
\bottomrule
\end{tabularx}
\caption{The claims of this paper, the evidence for each, and the data it covers. Hypotheses are defined in
Appendix~\ref{app:register}.}
\label{tab:claims}
\end{table}

\paragraph{Scope.}
This is not a new fitting method. Its contributions are a protocol (definitions, instruments and controls
for an identifiability claim about a fitted twin), a body of evidence on one model family and observable
under that protocol, and an honest account of what failed. The claims are limited to what the three
cohorts, the observation windows and the model family used here support (Section~\ref{sec:limits}).
